% !TeX root = ../机器学习笔记.tex
% 由主文档通过 \input 载入；请编译主文档。
\setcounter{chapter}{8} % 保留原书章号。
\chapter{贝叶斯决策与概率生成分类}
\label{chap:ml-generative}

概率生成分类先描述各类数据如何分布，再依据观测推断类别。模型是否贴近数据，决定了这种间接预测能否可靠。

\section{贝叶斯决策与分类代价}

设允许的行动为 $a$，真实类别为 $c$，损失为 $L(a,c)$。给定输入 $x$，行动的条件风险为 $r(a\mid x)=\sum_cL(a,c)P(Y=c\mid X=x)$。最优行动逐点最小化此风险；这是贝叶斯决策，不要求参数本身也使用贝叶斯方法估计。

零一损失下 $r(a\mid x)=1-P(Y=a\mid X=x)$，因此取最大后验类别。若允许“拒绝判断”，可为该行动设置固定代价，再与分类行动比较。概率本身描述相信什么，损失矩阵描述错误有多严重，决策同时依赖两者。

\section{生成方法与判别方法}

生成方法（generative method）建模 $p(x,y)=p(y)p(x\mid y)$，然后使用贝叶斯公式；判别方法（discriminative method）直接建模 $p(y\mid x)$ 或决策边界。生成模型可以描述某类输入的分布，但错误的输入分布假设也会影响后验；判别模型集中于预测目标，却不会自动提供完整联合分布。

估计的类别比例 $\widehat\pi_c=n_c/n$ 只是训练分布的先验估计。若使用场景的比例不同，即使类条件分布保持不变，后验仍会改变。两类方法没有不依赖样本量、模型设定和任务损失的统一优劣顺序。

\section{朴素贝叶斯及条件独立假设}

朴素贝叶斯（naive Bayes）假设给定类别后特征条件独立，于是
\[
 P(Y=c\mid x)\propto\pi_c\prod_{j=1}^d p_j(x_j\mid c),\qquad
 s_c(x)=\log\pi_c+\sum_j\log p_j(x_j\mid c).
\]
比较对数分数可避免长乘积下溢。Bernoulli 模型适合二元出现标记，Gaussian 模型为每类每个连续特征估计均值与方差，多项式模型则处理给定文档长度下的词频计数。不同观测机制对应不同似然，不能只因特征看起来是数字就混用。

对有 $V$ 个取值的类别型特征，类内取值 $v$ 的平滑估计为 $(N_{cv}+\alpha)/(N_c+V\alpha)$，$\alpha>0$ 防止未出现值得到零概率。条件独立不要求特征边缘独立；假设错误时，重复或高度相关的证据可能使概率过度集中，分类正确并不等于概率准确。

\section{高斯判别分析}

若 $X\mid Y=c\sim N(\mu_c,\Sigma)$ 且各类共享正定协方差，则删去与 $c$ 无关的项后，判别分数为
\begin{equation}
 \delta_c(x)=x^\top\Sigma^{-1}\mu_c
 -\tfrac12\mu_c^\top\Sigma^{-1}\mu_c+\log\pi_c.
 \label{eq:ml-gaussian-discriminant}
\end{equation}
两类分数之差对 $x$ 线性，因此称线性判别分析（linear discriminant analysis）。各类协方差不同则保留 $-\frac12\log|\Sigma_c|-\frac12(x-\mu_c)^\top\Sigma_c^{-1}(x-\mu_c)$，得到二次判别分析（quadratic discriminant analysis）。

类内样本数少于特征数时，经验协方差可能奇异；可使用共享协方差或向 $\tau I$ 收缩，但这改变了模型估计。协方差更灵活会增加参数数量，未必改善小样本预测。

\section{假设检验、诊断与应用限制}

生成模型的诊断包括检查类内相关性、尾部异常、零方差特征与类别覆盖。有限样本中“没有检测到违反假设”不等于假设成立，实际比较还应使用独立评价。概率校准可揭示仅看准确率无法发现的过度自信。

若仅类别先验改变且 $p_{\rm target}(x\mid c)=p_{\rm source}(x\mid c)$，可用目标先验替换原先验进行预测；若特征条件分布也改变，这种修正一般不足。训练中没有出现的类别，更不能仅靠修改先验获得可靠的类条件模型。

\section{例题、练习与学习检查}

\begin{example}[代价改变决策]
给定 $P(Y=1\mid x)=0.4$，正确分类代价为零，漏判代价为五，误报代价为一。
\end{example}
\begin{solve}
判正风险为 $1\times0.6=0.6$，判负风险为 $5\times0.4=2$，所以应判正，尽管正类后验小于二分之一。这不违反最大后验规则，因为后者只针对零一损失。
\end{solve}

习题：推导共享 Gaussian 协方差下的线性边界；对两个二元特征的计数表计算加一平滑后验；将一个特征复制两次，分析朴素贝叶斯如何重复计入证据，并比较分类结果与概率变化。





